
This study investigated cross-dimensional effects when applying targeted trustworthiness interventions to language models. Experiments across three transformer architectures (GPT-2, OPT, Pythia) tested whether LoRA fine-tuning targeting truthfulness produced predictable effects on fairness and robustness.

Results showed that representation changes propagated throughout network layers (100\% coverage, mean change 0.143), but cross-dimensional correlations varied by dimension pair. Truthfulness-fairness showed near-zero correlation (r = 0.034), truthfulness-robustness showed strong negative correlation in one architecture (r = -0.997, p = 0.051, not statistically significant), and fairness-robustness showed negative correlation in two of three architectures (replication rate 67\%, significant in one model).

These findings suggest cross-dimensional effects exist but are not uniform. However, interpretation is constrained by limited statistical power (n = 3-5 replicates), architecture coverage limited to transformers, single intervention method (LoRA), and evaluation restricted to three dimensions using benchmark proxies.

Future work could address these limitations through larger sample sizes, broader architecture coverage (including non-transformer models), comparison of intervention methods (LoRA vs. full fine-tuning vs. prompting), expanded dimension coverage, and evaluation in deployment contexts beyond benchmarks. Whether the observed patterns represent fundamental properties of trustworthiness trade-offs or artifacts of experimental design remains an open question requiring additional investigation with adequate statistical power.
